\documentclass[12pt]{article}
\usepackage{amsmath}
\usepackage{amsfonts}
\usepackage{amssymb}
\usepackage{geometry}
\usepackage{tikz}
\geometry{a4paper, margin=1in}
\pagenumbering{gobble}

\title{02YMECH - Homework 6}
\date{Assigned for the week of Oct 27, 2025}
\author{}
\begin{document}
\maketitle

\section*{Questions}
1 - Consider the following two diagrams:
\vspace{0.4cm}

\begin{minipage}{0.49\textwidth}
\centering
\pgfimage[width=\linewidth]{do.png}
\end{minipage}\hfill
\begin{minipage}{0.49\textwidth}
\centering
\pgfimage[width=\linewidth]{ip.png}
\end{minipage}

\vspace{0.5cm}
\noindent
Comment on the time evolution of the system in both the cases concisely. Which one would you classify as stable and which one as unstable? Justify your answer.

\vspace{0.5cm}
\noindent
2- A point mass moves in the plane with polar coordinates $(r,\theta)$ relative to a fixed origin. Let
\[
r(t)=r_0+v_r\,t,\qquad \theta(t)=\Omega\,t+\tfrac{1}{2}\alpha t^2,
\]
where \(r_0>0\), \(v_r\), \(\Omega\), and \(\alpha\) are constants. 

\begin{enumerate}
  \item[(a)] Write the general expressions for the velocity and acceleration in polar coordinates.
  \item[(b)] Find the radial and angular components of the velocity and acceleration as functions of time.
  \item[(c)] Compute the speed \(v(t)=\|\mathbf v\|\) and its time derivative \(\dot v(t)\). State the condition under which \(\dot v\) is positive, negative, or zero.

  \item[(d)] Find the time(s) \(t^\ast\) (if any) at which the radial component of acceleration \(a_r\) vanishes. State the condition on parameters for such a time to exist (with \(t^\ast\ge 0\)).
\end{enumerate}
\end{document}
